\section{Exact output: localized acceptance and polynomial boundary optima}\label{app:exact}

\subsection{Proof of Corollary~\ref{cor:local}}\label{app:localized}
We use the notation of Section~\ref{sec:localized} and
Corollary~\ref{cor:local}, and the bounds of Lemma~\ref{lem:commonmesh} for a
trial of CT (Algorithm~\ref{alg:ct}) with a common mesh with
$8\kappa\theta^2\le1$; the center of stage $j$ is $c=c^{(j)}$, and
$x^*\in X^{(j)}$.

\begin{proof}[Proof of Corollary~\ref{cor:local}]
\emph{Facts about $x^*$.} By Lemma~\ref{lem:growthcert}(b),
$\nabla_{J_0}F(x^*)=0$ and $H_{J_0J_0}\succeq2gI$; in particular
$H_{ii}\ge2g>0$ for $i\in J_0$, so $J_0\subseteq P$. Let $i\notin P$. If
$\ell_i<x^*_i<u_i$, the two points obtained from $x^*$ by setting coordinate
$i$ to $\ell_i$ and to $u_i$ are feasible and differ from $x^*$, so their
values exceed $\OPT$; but $F$ is
concave along coordinate $i$, so $\OPT=F(x^*)$ is at least the smaller of the
two values, a contradiction. So $x^*_i\in\{\ell_i,u_i\}$ for $i\notin P$. The
coordinates therefore split into $I_Z\cap P$, the coordinates $i\notin P$
(all at a bound), $J_0$ and $A$.

Fix a stage $j$ with $h_j\le h^*$ that ends by filtering, and put
$\omega_j=\sqrt{n\kappa}\,h_j\le\omega^*$. Then $\omega_j\le\frac12$ if
$I_Z\cap P\ne\emptyset$,
$\omega_j\le\delta_X/6$ if $I_C\cup(I_Z\setminus P)\ne\emptyset$, and
$\omega_j\le\lambda_A/(5\Gamma)$ if $A\ne\emptyset$. By (G2),
$\abs{(y_j)_i-x^*_i}\le\omega_j$ and $\abs{c_i-x^*_i}\le2\omega_j$ for every $i$, and
every node $v$ of $G_i$ with $m_i(v)\le U$ satisfies
$\abs{v-x^*_i}^2\le\frac{17}{15}\omega_j^2$, because the grid point that attains
$m_i(v)$ has $Q\le U$.

\emph{Coordinates $i\notin P$.} Here $G_i=\{\ell_i,u_i\}$ and $x^*_i$ is one
of the two nodes. The other node $v$ has
$\abs{v-x^*_i}=u_i-\ell_i\ge\delta_X\ge6\omega_j>\sqrt{17/15}\,\omega_j$, so
$m_i(v)>U$. The node $x^*_i$ has $m_i(x^*_i)\le F(x^*)=\OPT\le U$ by
Proposition~\ref{prop:cellwise}(c), because $x^*\in X^{(j)}$. Hence
$\tilde X_i=\{x^*_i\}$.

\emph{Coordinates $i\in I_Z\cap P$.} Every node $v\ne x^*_i$ has
$\abs{v-x^*_i}\ge1>\sqrt{17/15}\cdot\frac12\ge\sqrt{17/15}\,\omega_j$, so
$m_i(v)>U$. The center $c_i$ is an integer with $\abs{c_i-x^*_i}\le2\omega_j\le1$,
and $h_j\le\omega_j\le\frac12$. The grid steps
$\max\{1,\lfloor h_j+\theta t\rfloor\}$ at distance $t\le5$ from $c_i$ equal
one, because $h_j+\theta t\le\frac12+\frac54<2$. So every integer of
$X^{(j)}_i$ within distance $6$ of $c_i$ is a node. By (G4) and (G2),
$X^{(j+1)}_i$ lies within $1+4.2\omega_j+\omega_j\le3.6$ of $x^*_i$, hence within
$5$ of $c_i$. So every integer of $X^{(j+1)}_i$ is a node, and none lies
strictly inside a grid interval. Thus
$Z_i\cap X^{(j+1)}_i\subseteq\{x^*_i\}$, and $x^*_i$ belongs to it because
$x^*\in X^{(j+1)}$ and $m_i(x^*_i)\le\OPT\le U$. Hence $\tilde X_i=\{x^*_i\}$.

\emph{Continuous coordinates $i\in P$.} Here $\tilde X_i=X^{(j+1)}_i$, which
contains $x^*_i$ and, by (G4) and (G2), lies within
$4.2\omega_j+\omega_j<\delta_X$ of $x^*_i$. It has positive length, because
$X_i$ does and filtering replaces an interval of positive length by a hull of
grid intervals of positive length. So $i\in J_+$, and
$\tilde X_i$ contains a bound of $X_i$ exactly when $i\in A$, namely
the bound $x^*_i$. Hence $J_+=J_0\cup A$.

\emph{The face candidate.} Since $y_j\in\tilde X$, the face candidate has
$\tilde x_i=(y_j)_i=x^*_i$ for every integer coordinate and every $i\notin P$,
and $\tilde x_i=x^*_i$ for $i\in A$, while the free set is $J_f=J_0$.
The matrix $H_{J_0J_0}\succeq2gI$ is nonsingular, and $x^*$ satisfies
$\nabla_{J_0}F(x^*)=0$ and agrees with $\tilde x$ outside $J_0$, so
$\tilde x=x^*$.

\emph{The test.} We have $x^*\in\tilde X$ and $F(x^*)=\OPT\le U$. For
$i\in J_0$, $\zeta_i=\partial_iF(x^*)=0$, so case (i), (ii) or (iii) applies
with $\mu_i=0$. For $i\in A$, $x^*_i$ is an endpoint of $\tilde X_i$,
and since $x^*$ minimizes $F$ along coordinate $i$ in $X_i$, $\zeta_i\ge0$ at
$\ell_i$ and $\zeta_i\le0$ at $u_i$; so case (i) or (ii) applies. The interval
$\tilde X_i$ lies on one side of $x^*_i$, so its length is at most $5.2\omega_j$
and $\mu_i\ge\lambda_A/(5.2\omega_j)\ge\frac{5}{5.2}\Gamma\ge\Gamma/2$.

Let $N=H_{J_+J_+}+2\diag(\mu_{J_+})$. If $A=\emptyset$, then
$N=H_{J_0J_0}\succeq2gI$. If $J_0=\emptyset$, then
$N\succeq(2\min_{i\in A}\mu_i-\norm{H_{AA}}_2)I\succeq0$.
Otherwise $H_{J_0J_0}\succ0$, and $N\succeq0$ if the Schur complement
$H_{AA}+2\diag(\mu_{A})
-H_{AJ_0}H_{J_0J_0}^{-1}H_{J_0A}$ is positive semidefinite.
Its smallest eigenvalue is at least
$2\min_{i\in A}\mu_i-\norm{H_{AA}}_2
-\norm{H_{J_0A}}_2^2/(2g)=2\min_{i\in A}\mu_i-\Gamma\ge0$, using
$\norm{H_{J_0J_0}^{-1}}_2\le1/(2g)$. All hypotheses of
Proposition~\ref{prop:local} hold.

\emph{Stage count.} The least $j$ with $h_j=s2^{-j}\le h^*$ is at most
$\max\{0,\lceil\log_2(s/h^*)\rceil\}$. By Corollary~\ref{cor:height}(a), the
isolated minimizer $x^*$ and all endpoints lie in $\rho^{-1}\Z$ for an integer
$\rho\le R$ with $\Delta\mid\rho$. A nonzero difference $x^*_i-e$ therefore has
absolute value at least $1/\rho$, so $\delta_X\ge1/R$; and
$\partial_iF(x^*)=\sum_kH_{ik}x^*_k+b_i\in(\Delta\rho)^{-1}\Z$, so
$\lambda_A\ge1/(\Delta R)$. If $A\ne\emptyset$, then $L>0$, hence
$L\ge2/\Delta$ and $1/g\le\kappa/L\le\Delta\kappa/2$; with
$\norm{H}_2\le n\max_{ik}\abs{H_{ik}}$ this gives
$\log_2\Gamma\le\poly(I)+\log_2\kappa$. Since $\log_2s$, $\log_2R$ and
$\log_2\Delta$ are polynomial in $I$, the least such $j$ is at most
$\poly(I)+O(\log\kappa)$.
\end{proof}
